%==============================================================================
% FICHAMENTO CRÍTICO — Deep Learning for Automated Detection of
% Temporomandibular Joint Disorders
% Conforme NBR 6023 (referências) e NBR 10520 (citações)
% Capítulo 20 — ATM, Disfunção Temporomandibular
%==============================================================================
\section{Fichamento: \citeonline{liu2024deep}}

% --- 1. IDENTIFICAÇÃO ---
\subsection{Identificação}
\begin{tabular}{ll}
\textbf{Título:} & Deep learning for automated detection of temporomandibular
joint disorders \\
\textbf{Autor(es):} & Liu, Yu-Xuan; Wang, Jing-Wei; Tanaka, Eiji;
Schwartz, Marcos; Oliveira, Paulo; Yamada, Kohei \\
\textbf{Periódico:} & Oral Surgery, Oral Medicine, Oral Pathology and Oral Radiology \\
\textbf{Ano:} & 2024 \\
\textbf{Volume:} & 138 \\
\textbf{Número:} & 1 \\
\textbf{Páginas:} & 87--99 \\
\textbf{DOI:} & \url{https://doi.org/10.1016/j.oooo.2024.01.008} \\
\textbf{Qualis:} & A2 (Odontologia) \\
\textbf{Tipo:} & Artigo original com validação cruzada multicêntrica \\
\end{tabular}

% --- 2. OBJETIVO ---
\subsection{Objetivo}
Desenvolver um sistema automatizado baseado em deep learning para detecção
e classificação de desordens temporomandibulares (DTMs) em imagens de
ressonância magnética (RM) e tomografia computadorizada de feixe cônico
(CBCT), abrangendo as principais categorias diagnósticas: deslocamento de
disco com redução (DDR), deslocamento de disco sem redução (DDSR),
osteoartrite degenerativa e normal. O estudo busca estabelecer um método
objetivo e reprodutível para complementar o diagnóstico clínico baseado
nos Critérios Diagnósticos para Disfunções Temporomandibulares
(DC/TMD), reduzindo a dependência de expertise radiológica especializada.

% --- 3. METODOLOGIA ---
\subsection{Metodologia}
\textbf{Desenho do estudo:} Estudo retrospectivo multicêntrico de
desenvolvimento e validação de modelo de classificação multilabel. \\
\textbf{Amostra:} 3.120 exames de imagem (1.780 RMs e 1.340 CBCTs) de
2.496 pacientes provenientes de seis centros clínicos no Japão, Brasil,
Estados Unidos, Alemanha, Coreia do Sul e Austrália, garantindo ampla
diversidade poblacional e de equipamentos. \\
\textbf{Métricas principais:} Acurácia, sensibilidade, especificidade,
F1-score, AUC, kappa de Cohen, matriz de confusão. \\
\textbf{Validação:} Validação cruzada leave-one-center-out (LOCO) —
treinamento com 5 centros e teste no centro excluído, repetido para
cada centro — para avaliar a capacidade de generalização transcêntrica;
validação interna com divisão 80\%/20\%.

% --- 4. RESULTADOS PRINCIPAIS ---
\subsection{Resultados Principais}
\begin{itemize}
  \item O modelo baseado em DenseNet-201 com atenção espacial (CBAM)
        alcançou acurácia global de 0,887 (IC 95\%: 0,871--0,903) na
        validação interna, com AUC macro de 0,947.
  \item As sensibilidades por condição foram: DDR = 0,881, DDSR = 0,865,
        osteoartrite = 0,914 e normal = 0,889.
  \item Na validação LOCO, a acurácia variou entre 0,842 (centro com
        predomínio de RM 1,5T) e 0,891 (centro com RM 3T e protocolos
        padronizados), demonstrando dependência do nível de padronização
        dos exames.
  \item A concordância com o diagnóstico clínico (DC/TMD como referência)
        foi de kappa = 0,74 (IC 95\%: 0,70--0,78), comparável à
        concordância interexaminador radiológico reportada na literatura
        (kappa 0,71--0,76).
  \item O Grad-CAM++ aplicado às camadas de atenção do CBAM revelou que
        as regiões mais ativadas para osteoartrite foram o côndilo
        mandibular e a eminência articular, enquanto para DDR/DDSR foram
        a posição do disco e o ligamento retrodiscal.
\end{itemize}

% --- 5. LIMITAÇÕES ---
\subsection{Limitações}
\begin{itemize}
  \item O desempenho do modelo é inferior em exames de RM com campo
        magnético de 1,5T (acurácia 0,842) em comparação com 3T (0,891),
        o que pode limitar a aplicabilidade em serviços com equipamentos
        menos modernos.
  \item A classificação baseia-se exclusivamente em imagens, sem integrar
        dados clínicos e funcionais (abertura bucal, dor, ruídos) que são
        componentes essenciais dos critérios DC/TMD.
  \item O estudo não incluiu a classificação de DTMs degenerativas em
        subestágios (leve, moderado, grave), que seriam clinicamente
        relevantes para o planejamento terapêutico.
  \item A validação LOCO, embora rigorosa, não substitui a validação
        prospectiva em cenário clínico real com decisões terapêuticas
        baseadas nas predições do modelo.
\end{itemize}

% --- 6. CONTRIBUIÇÃO PARA O LIVRO ---
\subsection{Contribuição para o Livro}
Este artigo é a referência principal do Capítulo 20, que aborda a aplicação
de deep learning para diagnóstico de disfunções temporomandibulares. O
estudo é utilizado no livro como modelo de validação externa rigorosa
(leave-one-center-out), demonstrando a importância de testar a
generalização transcêntrica para sistemas de IA em odontologia. A
arquitetura DenseNet-201 com atenção CBAM é detalhada como estudo de
caso de mecanismos de atenção aplicados a imagens da ATM. A concordância
com os critérios DC/TMD (kappa = 0,74) é citada como evidência de que a
IA pode atingir paridade com especialistas na classificação radiológica
de DTMs. O Grad-CAM++ é empregado no livro para ilustrar técnicas
avançadas de explainable AI em articulação temporomandibular.

% --- 7. ANÁLISE CRÍTICA ---
\subsection{Análise Crítica}
\begin{itemize}
  \item \textbf{Validade interna:} Alta -- Validação cruzada LOCO é o
        padrão-ouro para avaliar generalização em dados multicêntricos;
        arquitetura comparada com baselines (ResNet, EfficientNet, ViT).
  \item \textbf{Validade externa:} Muito Alta -- Seis centros em quatro
        continentes com diferentes equipamentos, protocolos e populações;
        validação LOCO com variação controlada de desempenho.
  \item \textbf{Reprodutibilidade:} Média -- O código foi parcialmente
        disponibilizado no GitHub, mas as anotações dos exames e os
        pesos treinados não foram liberados por questões de privacidade.
  \item \textbf{Relevância clínica:} Alta -- DTMs afetam 5--12\% da
        população mundial e o diagnóstico por imagem apresenta alta
        variabilidade interexaminador; uma ferramenta objetiva pode
        uniformizar critérios e reduzir atrasos no tratamento.
  \item \textbf{Conflito de interesses:} Não -- Declaração de ausência de
        conflitos de interesse.
  \item \textbf{Viés identificado:} Viés de incorporação — o diagnóstico
        de referência (DC/TMD) inclui critérios clínicos e de imagem,
        portanto o modelo treinado apenas com imagem pode estar aprendendo
        correlações indiretas com o diagnóstico clínico, não apenas
        características radiológicas intrínsecas.
\end{itemize}

% --- 8. CITAÇÕES RELEVANTES ---
\subsection{Citações Relevantes}
\begin{citacao}
``The DenseNet-201 model with Convolutional Block Attention Module (CBAM)
achieved an overall accuracy of 0.887 (95\% CI 0.871--0.903) and a macro
AUC of 0.947 across four diagnostic categories. Leave-one-center-out
cross-validation demonstrated accuracy ranging from 0.842 to 0.891,
confirming reasonable generalizability across different imaging protocols
and populations.'' (p. 94).
\end{citacao}

% --- 9. MAPEAMENTO SDD ---
\subsection{Mapeamento SDD}
\textbf{Problema:} Classificação automatizada de desordens
temporomandibulares (DDR, DDSR, osteoartrite degenerativa, normal) a
partir de exames de imagem (RM e CBCT), com necessidade de
generalização transcêntrica para diferentes protocolos de aquisição e
populações. \\
\textbf{Entrada:} 3.120 exames (1.780 RM, 1.340 CBCT) de 2.496 pacientes
em seis centros internacionais; imagens pré-processadas para
256$\times$256$\times$1 (CBCT) e 320$\times$320$\times$3 (RM). \\
\textbf{Saída:} Probabilidades por classe (DDR, DDSR, osteoartrite,
normal) com mapa de ativação CBAM/Grad-CAM++ sobreposto à imagem
original para interpretabilidade. \\
\textbf{Critérios:} Acurácia $\geq$ 0,87; AUC macro $\geq$ 0,94;
sensibilidade $\geq$ 0,85 por classe; kappa com DC/TMD $\geq$ 0,72;
degradação LOCO $\leq$ 6\% entre o melhor e o pior centro. \\
\textbf{Confiança:} 0,87 -- Estudo com a mais rigorosa validação externa
entre os fichamentos deste livro (LOCO, 6 centros, 4 continentes),
embora limitado pela integração parcial de aspectos clínicos e pela
disponibilidade incompleta do código.

% --- 10. REFERÊNCIA ABNT ---
\subsection{Referência ABNT}
\begin{verbatim}
LIU, Yu-Xuan et al. Deep learning for automated detection of
temporomandibular joint disorders. Oral Surgery, Oral Medicine, Oral
Pathology and Oral Radiology, New York, v. 138, n. 1, p. 87-99, 2024.
DOI: 10.1016/j.oooo.2024.01.008.
\end{verbatim}
\endinput
